% Part: methods
% Chapter: proofs
% Section: starting-proofs

\documentclass[../../../include/open-logic-section]{subfiles}

\begin{document}

\olfileid[hi]{mth}{prf}{str}

\olsection{प्रमाण आरंभ करना}

किंतु आरंभ कहाँ से करें?

आपको कुछ सिद्ध करने के लिए दिया गया है, इसलिए प्रमाण में उसका उल्लेख सबसे
अंत में होना चाहिए (निस्संदेह आरंभ में आप \emph{घोषित} कर सकते हैं कि उसे
सिद्ध करने जा रहे हैं, किंतु उसे मान्यता के रूप में उपयोग नहीं करना चाहते)।
कागज़ के नए पन्ने के नीचे वह लिखें जिसे सिद्ध करने का प्रयास कर रहे हैं---इस
प्रकार आपका लक्ष्य दृष्टि से ओझल नहीं होगा।

इसके बाद आपके पास उपयोग योग्य कुछ मान्यताएँ हो सकती हैं (अगले अनुभाग में
आपके किए जा रहे प्रमाण के \emph{प्रकार} पर चर्चा के समय यह और स्पष्ट होगा)।
उन्हें पृष्ठ के ऊपर लिखें और यह स्पष्ट करें कि वे मान्यताएँ हैं (अर्थात् यदि
$p$ मान रहे हैं, तो ``$p$ मानें'' या ``मान लें कि $p$'' लिखें)। अंततः प्रश्न
में कुछ ऐसी परिभाषाएँ हो सकती हैं जिन्हें जानना आवश्यक है। आपको किसी विशिष्ट
परिभाषा का उपयोग करने के लिए कहा जा सकता है, या जिन मान्यताओं अथवा निष्कर्ष
की ओर काम कर रहे हैं उनमें विभिन्न परिभाषाएँ हो सकती हैं। \emph{इन्हें
लिखें और सुनिश्चित करें कि आप इनका अर्थ समझते हैं।}

प्रमाण को कैसे व्यवस्थित करें, यह प्रश्न के रूप पर भी निर्भर होगा। अगला
अनुभाग वाक्य के प्रकार के आधार पर प्रमाण व्यवस्थित करने का विवरण देता है।

\end{document}
